
\begin{ledgroupsized}[r]{120mm}
\footnotesize
\pstart
\noindent\textbf{Überlieferung:}
\pend
\end{ledgroupsized}
\begin{ledgroupsized}[r]{114mm}
\footnotesize
\pstart
\parindent -6mm
\makebox[6mm][l]{\textit{E}}%
Erstdruck:
\cite{01023}\textit{AE} März 1694, S.\;112\textendash113.
Wie zu Leibnizens Zeiten üblich erschien die Rezension anonym.
\pend
\end{ledgroupsized}
\selectlanguage{latin}
\vspace{8mm}
\pstart%
\centering
\lbrack S.\;112\rbrack\
\hfill
\edtext{\cite{02001}\textit{Investigatio Momentorum, quibus gravia tendunt deorsum}, 
\hfill 
\phantom{\lbrack S.\;112\rbrack}
\linebreak
autore \indexN{\textls{Vanni}, Giovanni Francesco 1638\textendash1709}%
Jo.\;Franc. Vannio e Soc.\;Jesu, etc.
\linebreak
Romae apud D.\,A.\,Herculem, 1693. in 8. Plag.~7.}{%
\lemma{\textit{Investigatio} \lbrack...\rbrack\ Plag.~7}%
\Cfootnote{\indexN{\textls{Vanni}, Giovanni Francesco 1638\textendash1709}\textsc{G.\,F. Vanni}, \cite{02001}\textit{Investigatio momentorum quibus gravia tendunt deorsum}, Rom 1693.}}
\pend
\vspace{0.5\baselineskip}
\pstart\noindent
\edtext{\indexN{\textls{Vanni}, Giovanni Francesco 1638\textendash1709}%
Autor Clarissimus \cite{02001}librum suum Illustrissimo Viro, 
\indexN{\textls{Ciampini}, Giovanni Giusto 1633\textendash1698}\textls{Johanni Ciampino} dedicans,}{%
\lemma{Autor \lbrack...\rbrack\ dedicans}%
\Cfootnote{\cite{02001}a.\,a.\,O., S.\;7.}}
\edtext{inter caetera ejus in rem literariam et publicam merita memorat,
quod 
\cite{02101}\textit{Ephemerides Literatorum}
primus Romae\indexO{Rom} publici juris fecerit;}{%
\lemma{inter \lbrack...\rbrack\ fecerit}%
\Cfootnote{\cite{02001}a.\,a.\,O., S.\;7\,f.
Nach einer Spaltung in der Redaktion des ersten \cite{02101}\textit{Giornale de' Letterati} (Rom 1668--1675) gründete \indexN{\textls{Ciampini}, Giovanni Giusto 1633\textendash1698}Ciampini eine neue Reihe des \textit{Giornale},  in der u.\,a. \indexN{\textls{Vanni}, Giovanni Francesco 1638\textendash1709}Vanni und \indexN{\textls{Eschinardi}, Francesco 1623\textendash1703}Eschinardi publizierten.}}
\edtext{et quod mox Praesulis dignitate
ac pluribus muneribus auctus Academiam apud
\indexS{Collegium Urbanum de propaganda fide}%
Collegium Urbanum de propaganda fide,
in qua de rebus ecclesiasticis\indexS{res ecclesiastica} accuratissime disseratur, 
alteramque in aedibus suis, quae Physico-Mathematica\indexS{Academia Physico-Mathematica} 
nominetur, condiderit.}{%
\applabel{FnC:AE_1694_112-113_ciampini}%
\lemma{et quod \lbrack...\rbrack\  condiderit}%
\Cfootnote{\cite{02001}a.\,a.\,O., S.\;8.
Die "Accademia Fisicomatematica" versammelte sich von August 1677 bis 1689 bei Ciampini, dann bei Kardinal \indexN{\textls{Ottoboni}, Pietro 1667\textendash1740}Ottoboni, und löste sich nach Ciampinis Tod um 1698 auf (siehe \textsc{Robinet}, \textit{Iter Italicum}, Florenz 1988, S.\;42--121.)}}
Additque 
\edtext{in hac suam hanc de momentis\indexS{momentum} philosophiam audientia
\lbrack S.\;113\rbrack\ facta vitali aura ac spiritu frui coepisse.}{%
\lemma{in hac \lbrack...\rbrack\  coepisse}%
\Cfootnote{\cite{02001}a.\,a.\,O., S.\;9.}}
\pend
\pstart
Caeterum quod materiam \cite{02001}libri attinet, 
\edtext{notum est dudum vel ex \cite{01023}\textit{Actis} nostris, 
\indexN{\textls{Vanni}, Giovanni Francesco 1638\textendash1709}%
Autorem scriptis quibusdam editis impugnasse 
receptam doctrinam\indexS{doctrina recepta de momentis gravium in planis inclinatis} de momentis gravium\indexS{momentum gravis} praesertim in planis inclinatis,\indexS{planum inclinatum}}{%
\lemma{notum \lbrack...\rbrack\ inclinatis}%
\Cfootnote{\indexN{\textls{Vanni}, Giovanni Francesco 1638\textendash1709}Vannis anonym erschienenem Buch 
\cite{02000}\textit{Exegeses physico-mathematicae}, Rom 1685, war die Flugschrift 
\cite{01995}\textit{Specimen libri de momentis gravium}, Rom 1684, vorangegangen.
Diese hatten die 
\cite{01023}\textit{Acta Eruditorum} (November 1684, S.\;511\textendash513) nebst einer ebenfalls anonymen kritischen Anmerkung von Leibniz (ebd., S.\;514) abgedruckt. Die Anmerkung erscheint in diesem Band als N.\;\ref{intro_vanni-specimen}, Leibnizens anonyme Rezension der \textit{Exegeses} als N.\;\ref{41199}.}}
ingeniosisque in eam instantiis usum.
Adversarios vero nactus est aliquot Viros doctos,
qui justis scriptis doctrinam communem\indexS{doctrina communis de momentis gravium} defenderunt,
atque illustrare ac promovere sunt nisi, quibus hic respondetur;
nimirum Viris clarissimis, Professoribus Matheseos publicis,
\edtext{\indexN{\textls{Giordani} (Giordano), Vitale 1633\textendash1711}Vitali Jordano Romae,}{%
\lemma{Vitali Jordano Romae}%
\Cfootnote{Siehe \indexN{\textls{Giordani} (Giordano), Vitale 1633\textendash1711}\textsc{V.~Giordani},
\cite{02110}\textit{De componendis gravium momentis dissertatio}, Rom 1687 und das umfangreichere
\cite{02109}\textit{Fundamentum doctrinae motus gravium}, Rom 1689.
Leibniz hat aus beiden Werken exzerpiert: siehe N.\;\ref{58253} in diesem Band.}}
\edtext{\indexN{\textls{Marchetti}, Alessandro 1633\textendash1714}Alexandro Marchetto Pisis,}{%
\lemma{Alexandro Marchetto Pisis}%
\Cfootnote{%
Siehe 
die Flugschrift \cite{02113}\textit{Conclusioni intorno a' momenti de' gravi sopra i piani declivi}, Florenz 1687
und das Buch \cite{02114}\textit{Prove delle conclusioni intorno a' momenti de' gravi sopra i piani declivi}, Florenz 1688. 
Beide Werke stammen eigentlich von \indexN{\textls{Marchetti}, Alessandro 1633\textendash1714}Alessandro Marchetti, wurden aber unter dem Decknamen seines dreizehnjährigen Sohnes \indexN{\textls{Marchetti}, Angelo 1674\textendash1753}Angelo veröffentlicht.
Siehe die Erläuterung zu N.\;\ref{41823}, \refnote{FnC:35_10_13_010-011_marchettus}.}}
et 
\edtext{\indexN{\textls{Spoleti}, Francesco, gest. 1712}Francisco Spoleto Patavii.}{%
\lemma{Francisco Spoleto Patavii}%
\Cfootnote{%
\indexN{\textls{Spoleti}, Francesco, gest. 1712}Spoleti reagierte bereits 1684 auf 
\indexN{\textls{Vanni}, Giovanni Francesco 1638\textendash1709}Vannis 
Thesen mit dem Opusculum \cite{02096}\textit{De momento gravis in plano inclinato}, Rom 1684, das Vanni in der \cite{02001}\textit{Investigatio} wiederum zitiert,
sowie der Dissertation \cite{02095}\textit{De momento quo gravia nituntur deorsum per datum planum}, Venedig 1686, die Gegenstand einer nicht veröffentlichen Besprechung durch Leibniz war (N.\;\ref{41200} in diesem Band).}}
Et 
\indexN{\textls{Giordani} (Giordano), Vitale 1633\textendash1711}Jordano quidem ac 
\indexN{\textls{Marchetti}, Alessandro 1633\textendash1714}Marchetto passim, 
\edtext{\indexN{\textls{Spoleti}, Francesco, gest. 1712}Spoleto vero p.~91.}{%
\lemma{Spoleto vero p.~91}%
\Cfootnote{\indexN{\textls{Vanni}, Giovanni Francesco 1638\textendash1709}\textsc{Vanni}, \cite{02001}\textit{Investigatio}, S.\;91\,f.}}
\pend
\pstart
Quin et 
\edtext{p.~94 respondetur}{%
\lemma{p.~94 respondetur}%
\Cfootnote{\cite{02001}a.\,a.\,O., S.\;94\,f.}}
iis, 
\edtext{quae R.\,P. 
\indexN{\textls{Kochański}, Adam Adamandus 1631\textendash1700}Adamus Adamandus Kochanskius 
in nostris \cite{01023}\textit{Actis} objecerat\lbrack;\rbrack}{%
\lemma{quae \lbrack...\rbrack\ objecerat}%
\Cfootnote{\indexN{\textls{Kochański}, Adam Adamandus 1631\textendash1700}\textsc{A.\,A. Kochański}, 
\cite{01996}\textit{Consideratio speciminis libri de momentis gravium, Authore J.\,F.\,V. Lucensi}, \cite{01023}\textit{AE} Juni 1685, S.\;262\textendash265.}}
passim etiam disputatur contra 
\indexN{\textls{Galilei} (Galilaeus, Galileus), Galileo 1564\textendash1642}Galilaeum, et speciatim 
\edtext{p.~108 contra quaedam ejus dicta in libro \cite{02085}\textit{de insidentibus humido}.}{%
\lemma{p.~108 \lbrack...\rbrack\ \textit{humido}}%
\Cfootnote{\indexN{\textls{Vanni}, Giovanni Francesco 1638\textendash1709}\textsc{Vanni}, \cite{02001}\textit{Investigatio}, S.\;108\textendash110. Siehe 
\indexN{\textls{Galilei} (Galilaeus, Galileus), Galileo 1564\textendash1642}\textsc{G.~Galilei},
\cite{02085}\textit{Discorso intorno alle cose che stanno in su l'acqua, o che in quella si muovono}, Florenz 1612 (\cite{00048}\textit{GO}~IV, S.\;57\textendash141).}}
Sed et 
\edtext{a \indexN{\textls{Stevin} (Stevinus), Simon 1548\textendash1620}Stevini p.~105,}{%
\lemma{a Stevini p.~105}%
\Cfootnote{\indexN{\textls{Vanni}, Giovanni Francesco 1638\textendash1709}\textsc{Vanni}, \cite{02001}\textit{Investigatio}, S.\;104\,f.; 
siehe
\indexN{\textls{Stevin} (Stevinus), Simon 1548\textendash1620}%
\textsc{S.~Stevin}, \cite{00196}\textit{De Beghinselen der Weeghconst}, Leiden 1586, 
bes. Buch~1, Teil~2, Theor.~XI Prop.~XIX mit Koroll.~I, S.\;40\textendash42
(\cite{02084}\textit{SPW} I, S.\;174\textendash179).}}
et 
\edtext{\indexN{\textls{Torricelli} (Torricellius), Evangelista 1608\textendash1647}Torricelli 
p.~72 argumentis se expedire conatur;}{%
\lemma{Torricelli \lbrack...\rbrack\  conatur}%
\applabel{FnC:AE_1694_112-113_Torricelli}%
\Cfootnote{\indexN{\textls{Vanni}, Giovanni Francesco 1638\textendash1709}\textsc{Vanni}, \cite{02001}\textit{Investigatio}, prop.~XXVIII mit scholion, S.\;70\textendash72. 
Siehe 
\indexN{\textls{Torricelli} (Torricellius), Evangelista 1608\textendash1647}\textsc{E.~Torricelli}, 
\textit{De motu gravium}, lib.~I (\cite{00304}\textit{Opera geometrica}, Florenz 1644, S.\;99;
\cite{00108}\textit{TO}~II, S.\;105).}}
quam feliciter, nos nobis judicium non sumimus.
\pend
\pstart
Ex propria etiam Societate\indexS{Societas Jesu} 
(ne \indexN{\textls{Kochański}, Adam Adamandus 1631\textendash1700}Adamandum repetamus)
\edtext{impugnat \indexN{\textls{Zucchi}, Niccolò 1586\textendash1670}Zuchium p.~40,}{%
\lemma{impugnat Zuchium p.~40}%
\Cfootnote{\indexN{\textls{Vanni}, Giovanni Francesco 1638\textendash1709}\textsc{Vanni}, \cite{02001}\textit{Investigatio}, eigentlich auf S.\;42. 
Siehe
\indexN{\textls{Zucchi}, Niccolò 1586\textendash1670}\textsc{N.~Zucchi},
\cite{02104}\textit{Nova de machinis philosophia}, Rom 1649, bes. pars~III, sec.~II, S.\;68\textendash80.}} %
\edtext{\indexN{\textls{Eschinardi}, Francesco 1623\textendash1703}Eschinardum p.~42 et alias,}{%
\lemma{Eschinardum p.~42 et alias}%
\applabel{FnC:AE_1694_112-113_Eschinardi}%
\Cfootnote{\indexN{\textls{Vanni}, Giovanni Francesco 1638\textendash1709}\textsc{Vanni}, \cite{02001}\textit{Investigatio}, eigentlich auf S.\;43\,f.
Siehe
\indexN{\textls{Eschinardi}, Francesco 1623\textendash1703}\textsc{F.~Eschinardi},
\textit{Cursus physicomathematicus}, tom.~I,
Rom 1689, bes. app.~IV, S.\;358\textendash361.}}
\edtext{\indexN{\textls{Casati}, Paolo 1617\textendash1707}Casatum p.~99.}{%
\lemma{Casatum p.~99}%
\Cfootnote{\indexN{\textls{Vanni}, Giovanni Francesco 1638\textendash1709}\textsc{Vanni}, \cite{02001}\textit{Investigatio}, S.\;99.
Siehe
\indexN{\textls{Casati}, Paolo 1617\textendash1707}\textsc{P.~Casati}, 
\cite{02106}\textit{Mechanicorum libri octo}, Lyon 1684, bes. lib.~I, cap.~XIII, S.\;79\textendash86.}}
Phrasibus sibi peculiaribus subinde utitur, 
videturque quaedam profunde in hoc argumento meditatus, ut optandum sit ad scientiae pomoeria
proferenda ingenium ejus porro converti.
\edtext{\indexN{\textls{Pappos} (Pappus) von Alexandreia um 300}Pappi
propositionem 9. libri octavi culpat\lbrack;\rbrack\
eamque,
licet a \indexN{\textls{Del Monte}, Guidobaldo 1545\textendash1607}Guidone Ubaldo de Monte 
in \cite{01009}\textit{Mechan.} prop.~2. de cochlea admissam,
tamen a \indexN{\textls{Cabeo}, Niccolò 1586\textendash1650}P.\,Nic. Cabeo
lib.~4. \cite{02107}\textit{Philos. Magneticae} c.~20 ad absurdum fuisse deductam;}{%
\lemma{Pappi \lbrack...\rbrack\ deductam}%
\Cfootnote{\indexN{\textls{Vanni}, Giovanni Francesco 1638\textendash1709}\textsc{Vanni}, \cite{02001}\textit{Investigatio}, \textit{Praefatio}, S.\;4.
Siehe
\indexN{\textls{Pappos} (Pappus) von Alexandreia um 300}\textsc{Pappos}, 
\cite{01015}\textit{Collectio} VIII, prop.~9 (Hultsch 1054\textendash1059);
\indexN{\textls{Del Monte}, Guidobaldo 1545\textendash1607}\textsc{G. Del~Monte}, 
\cite{01009}\textit{Mechanicorum liber}, Pesaro 1577, \textit{De cochlea}, prop.~II, Bl.~124~v\textsuperscript{o};
\indexN{\textls{Cabeo}, Niccolò 1586\textendash1650}\textsc{N.~Cabeo}, 
\cite{02107}\textit{Philosophia magnetica}, Ferrara 1629, lib.~IV, cap.~XX, S.\;338\textendash346.}}
\edtext{quin et assumtum \indexN{\textls{Pappos} (Pappus) von Alexandreia um 300}Pappi a 
\indexN{\textls{Galilei} (Galilaeus, Galileus), Galileo 1564\textendash1642}Galilaeo 
reprehendi in lib. \cite{02108}\textit{de Scientia Mechanica} monet.}{%
\lemma{quin \lbrack...\rbrack\ monet}%
\Cfootnote{\indexN{\textls{Vanni}, Giovanni Francesco 1638\textendash1709}\textsc{Vanni}, \cite{02001}\textit{Investigatio}, S.\;4\,f.
Siehe \indexN{\textls{Galilei} (Galilaeus, Galileus), Galileo 1564\textendash1642}\textsc{G.~Galilei},
\cite{02108}\textit{Della scienza mecanica}, Bologna 1649, S.\;25 (\textit{GO}~II, S.\;181).}}
Caeterum 
\edtext{%
proportionem momentorum,\indexS{momentum} quae gravia exercent in planis inclinatis,\indexS{planum inclinatum}\indexS{proportio momentorum gravium in planis inclinatis}
cum \indexN{\textls{Cardano} (Cardanus), Girolamo (Geronimo) 1501\textendash1576}Cardano
metitur angulis planorum ad horizontem,\indexS{horizon}\indexS{angulus plani ad horizontem}
non vero angulorum sinubus\indexS{sinus anguli} cum 
\indexN{\textls{Tartaglia} (Tartalea), Niccolò 1499 o. 1500\textendash1557}Tartalea,
cui jam praeiverat 
\indexN{\textls{Nemorarius}, Jordanus 13. Jh.}Jordanus 
suo aevo celebris Mathematicus,
\indexN{\textls{Stevin} (Stevinus), Simon 1548\textendash1620}Stevinusque ac 
\indexN{\textls{Galilei} (Galilaeus, Galileus), Galileo 1564\textendash1642}Galilaeus
et plerique omnes sunt secuti.}{%
\lemma{proportionem \lbrack...\rbrack\ secuti}%
\Cfootnote{\indexN{\textls{Vanni}, Giovanni Francesco 1638\textendash1709}\textsc{Vanni}, \cite{02001}\textit{Investigatio}, S.\;5\,f.
Siehe
\indexN{\textls{Cardano} (Cardanus), Girolamo (Geronimo) 1501\textendash1576}\textsc{G.~Cardano}, 
\cite{02125}\textit{Opus novum de proportionibus}, prop.~LXXII, S.\;63;
\indexN{\textls{Tartaglia} (Tartalea), Niccolò 1499 o. 1500\textendash1557}\textsc{N.~Tartaglia}, %
\cite{02016}\textit{Quesiti et inventioni diverse}, 2.~Ausgabe, Venedig 1554,
Libro Ottavo, Quesito XLII Propositione XV, 
Bl.~97~r\textsuperscript{o}\textendash97~v\textsuperscript{o};
\indexN{\textls{Nemorarius}, Jordanus 13. Jh.}\textsc{Jordanus Nemorarius}, \cite{02128}\textit{De ratione ponderis}, 1.~Hälfte 13.~Jh., "quinta suppositio" (in Tartaglias Edition \textit{Iordani opusculum de ponderositate}, Venedig 1565, auf S.\;3)
\indexN{\textls{Stevin} (Stevinus), Simon 1548\textendash1620}%
\textsc{Stevin}, \cite{00196}\textit{De Beghinselen der Weeghconst}, S.\;40\textendash42;
\indexN{\textls{Galilei} (Galilaeus, Galileus), Galileo 1564\textendash1642}\textsc{Galilei},
\cite{00050}\textit{Discorsi}, Giornata Terza, De motu naturaliter accelerato; 
\glqq Aggiunta postuma dell'Autore\grqq\ nach dem Scholion zum 2.~Theorem (der Zusatz erschien erst in der posthumen Werkausgabe: \cite{01084}\textit{Opere}, Bologna 1656, Bd.~II, S.\;132\lbrack ter\rbrack\textendash134; \cite{00048}\textit{GO} VIII, S.\;216\textendash219 im Apparat).}}
\pend
\pstart
\edtext{Argumentum in his \cite{01023}\textit{Actis} oppositum, (quo etiam 
\indexN{\textls{Torricelli} (Torricellius), Evangelista 1608\textendash1647}Torricellius 
jam usus hic refertur) quod gravia\indexS{gravia in aequilibrio} tunc censenda sint esse
in aequilibrio,\indexS{aequilibrium} cum commune eorum centrum gravitatis\indexS{centrum gravitatis commune gravium} descendere non potest,\indexS{centrum gravitatis descendere nequiens}
eo effugere satagit \indexN{\textls{Vanni}, Giovanni Francesco 1638\textendash1709}Vannius,
quod duo gravia, filo connexa\indexS{gravia duo filo connexa} et in duobus diversis planis
descendere conantia,\indexS{gravia in diversis planis descendere conantia} 
communi gravitatis centro (ut ipse arbitratur) careant.\indexS{centrum gravitatis commune}\indexS{gravia centro gravitatis commune carentia}}{%
\lemma{Argumentum \lbrack...\rbrack\ careant}%
\applabel{FnC:AE_1694_112-113_argumentum}%
\Cfootnote{Leibniz spielt auf seine anonym erschienene \textit{Demonstratio geometrica regulae apud Staticos receptae} an (\textit{AE} November 1685, S.\;501\textendash505, N.\;\ref{49409} in diesem Band, bes. \refpassage{AE_1685_501-505_15a}{AE_1685_501-505_15b}.
Allerdings erwähnt \indexN{\textls{Vanni}, Giovanni Francesco 1638\textendash1709}Vanni 
in diesem Zusammenhang keine Publikation aus den \cite{01023}\textit{Acta Eruditorum}, sondern \indexN{\textls{Torricelli} (Torricellius), Evangelista 1608\textendash1647}Torricellis \textit{De motu gravium}: siehe die Erläuterung zu \refnote{FnC:AE_1694_112-113_Torricelli}.}}
\pend
